%==========================================================
%  CHAPITRE 09 — CONCLUSION GÉNÉRALE ET PERSPECTIVES
%==========================================================
\chapter{Conclusion générale et perspectives}
\label{chap:conclusion}

%──────────────────────────────────────────────────────────
\section{Bilan du projet}
\label{sec:ccl:bilan}
%──────────────────────────────────────────────────────────

Ce stage de \todo{durée} au sein de \textbf{\CompanyName} avait pour
objectif principal \todo{rappeler l'objectif central en une phrase}.

Au terme de cette expérience, les principaux résultats obtenus sont :

\begin{itemize}
  \item \todo{Résultat 1 : livrable principal et son impact mesurable}
  \item \todo{Résultat 2 : amélioration quantifiée (temps, coût, qualité)}
  \item \todo{Résultat 3 : contribution à l'équipe ou au système existant}
  \item \todo{Résultat 4 : apprentissage ou montée en compétence notable}
\end{itemize}

\todo{Rédiger un paragraphe de synthèse (5--8 phrases) qui évalue
l'atteinte des objectifs fixés en introduction, en soulignant les
écarts éventuels et leurs causes.}

%──────────────────────────────────────────────────────────
\section{Apports personnels}
\label{sec:ccl:apports}
%──────────────────────────────────────────────────────────

\subsection{Compétences techniques acquises}

\todo{Dresser la liste des compétences techniques nouvelles ou consolidées :
langages, frameworks, outils, pratiques de développement (TDD, CI/CD,
clean code, etc.).}

\begin{itemize}
  \item \textbf{Maîtrise approfondie de} \tech{\todo{framework / langage}}
  \item \textbf{Découverte et mise en pratique de} \tech{\todo{outil / pratique}}
  \item \textbf{Compréhension de} \todo{concept architectural ou méthodologique}
\end{itemize}

\subsection{Compétences transversales}

\todo{Mentionner les soft-skills développées : communication, gestion du
temps, travail en équipe, autonomie, animation de réunions, etc.}

\begin{itemize}
  \item \textbf{Communication :} participation aux daily standups et
        présentations d'avancement devant l'équipe.
  \item \textbf{Gestion de projet :} utilisation de \todo{Jira/Trello}
        pour le suivi des tâches, respect des délais de sprint.
  \item \textbf{Travail en équipe :} collaboration en mode pair-programming
        avec \todo{collègue}, revues de code croisées.
  \item \todo{Autre compétence transversale}
\end{itemize}

%──────────────────────────────────────────────────────────
\section{Difficultés rencontrées et leçons apprises}
\label{sec:ccl:difficultes}
%──────────────────────────────────────────────────────────

\todo{Synthétiser les principales difficultés (techniques, organisationnelles,
humaines) et les solutions trouvées. Se montrer critique et honnête :
ce qui a bien fonctionné, ce qui aurait pu être fait différemment.}

\begin{itemize}
  \item \textbf{Difficulté technique :} \todo{description + enseignement tiré}
  \item \textbf{Difficulté méthodologique :} \todo{description + enseignement tiré}
  \item \textbf{Difficulté organisationnelle :} \todo{description + enseignement tiré}
\end{itemize}

%──────────────────────────────────────────────────────────
\section{Perspectives et améliorations futures}
\label{sec:ccl:perspectives}
%──────────────────────────────────────────────────────────

Bien que les objectifs principaux aient été atteints, plusieurs pistes
d'amélioration et d'évolution sont envisagées à court et moyen terme :

\begin{enumerate}
  \item \textbf{Court terme (0--3 mois) :}
    \begin{itemize}
      \item \todo{Fonctionnalité non implémentée par manque de temps}
      \item \todo{Amélioration de performances identifiée en tests de charge}
      \item \todo{Correction de bugs mineurs détectés en recette}
    \end{itemize}

  \item \textbf{Moyen terme (3--12 mois) :}
    \begin{itemize}
      \item \todo{Extension du périmètre fonctionnel — nouvelle feature}
      \item \todo{Migration vers une architecture plus scalable}
      \item \todo{Intégration avec un système tiers}
    \end{itemize}

  \item \textbf{Long terme (> 1 an) :}
    \begin{itemize}
      \item \todo{Vision stratégique : digitalisation complète du processus X}
      \item \todo{Adoption de l'IA / ML pour automatiser la prise de décision}
      \item \todo{Autre perspective}
    \end{itemize}
\end{enumerate}

%──────────────────────────────────────────────────────────
\section{Mot de clôture}
\label{sec:ccl:mot}
%──────────────────────────────────────────────────────────

\begin{tcolorbox}[
  enhanced,
  colback=LightBg, colframe=ChapBlue,
  fontupper=\itshape,
  width=\textwidth
]
\todo{Rédiger un court paragraphe conclusif personnel (3--5 phrases)
qui exprime votre ressenti global, l'impact de ce stage sur votre
projet professionnel, et vos ambitions pour la suite de votre parcours.
Exemple : « Ce stage a été une expérience fondatrice qui a confirmé
mon attrait pour [domaine]. Il m'a permis de \ldots\ et m'a donné
envie de \ldots »}
\end{tcolorbox}

\clearpage
